\subsection{QUADRATIC ALGEBRAS}

Let \(\alpha, \beta\) be two elements of A and \((e_1, e_2)\) the canonical basis of \(\mathbf{A}^2\) . The quadratic algebra of type \((\alpha, \beta)\) over A is the A-module \(\mathbf{A}^2\) with the algebra structure defined by the multiplication table (§ 1, no. 7)

\[
e _ {1} ^ {2} = e _ {1}, \qquad e _ {1} e _ {2} = e _ {2} e _ {1} = e _ {2}, \qquad e _ {2} ^ {2} = \alpha e _ {1} + \beta e _ {2}.\tag{2}
\]

An A-algebra E isomorphic to a quadratic algebra is also called a quadratic algebra. It amounts to the same to say that E admits a basis of two elements one of which is the unit element.

It can be shown that every unital A-algebra which admits a basis of two elements is a quadratic algebra (Exercise 1).

If a basis \((e_{1}, e_{2})\) of an A-algebra has multiplication table (2), it is called a basis of type \((\alpha, \beta)\) . By an abuse of language, a quadratic algebra is said to be of type \((\alpha, \beta)\) when it has a basis of type \((\alpha, \beta)\) .

PROPOSITION 1. A quadratic algebra E is associative and commutative.

The fact that E is commutative follows from the equation \(e_{1}e_{2}=e_{2}e_{1}\) in (2); similarly, to verify associativity, it suffices to see that \(x(yz)=(xy)z\) when x,y,z are each equal to \(e_{1}\) or \(e_{2}\) . Now, this relation is obvious if at least one of the elements x,y,z is equal to \(e_{1}\) ; it is also true for x=y=z=e₂ since E is commutative; whence the proposition.

Let \(e\) denote the unit element of a quadratic algebra \(\mathbf{E}\) and let \((e, i)\) be a basis of \(\mathbf{E}\) of type \((\alpha, \beta)\) ; every other basis of \(\mathbf{E}\) containing \(e\) is therefore of the form \((e, j)\) with \(j = \gamma e + \delta i\) (II, §7, no. 2, Corollary to Proposition 3); moreover, for \((e, j)\) to be a basis of \(\mathbf{E}\) , it is necessary and sufficient that \(\delta\) be invertible in \(\mathbf{A}\) ; the condition is obviously sufficient; conversely, if \(i\) is the canonical image of \(i\) in \(\mathbf{E}/\mathbf{A}e\) , \(i\) and \(\bar{j} = \delta i\) must each form a basis of \(\mathbf{E}/\mathbf{A}e\) , whence the necessity of the condition. Then

\[
j ^ {2} = (\gamma^ {2} + \alpha \delta^ {2}) e + (2 \gamma \delta + \beta \delta^ {2}) i = (\alpha \delta^ {2} - \gamma^ {2} - \beta \gamma \delta) e + (2 \gamma + \beta \delta) j;
\]

thus it is seen that E is of type

\[
(\alpha \delta^ {2} - \gamma^ {2} - \beta \gamma \delta , 2 \gamma + \beta \delta)\tag{3}
\]

for all invertible \(\delta \in \mathbf{A}\) and all \(\gamma \in \mathbf{A}\) . In particular, if \(\mathbf{E}\) is of type \((\alpha, 2\beta')\) , i is also of type \((\alpha + \beta'^2, 0)\) as is seen by taking \(\gamma = -\beta'\) and \(\delta = 1\) .

PROPOSITION 2. Let E be a quadratic A-algebra and e its unit element. For all \(u \in \mathbf{E}\) , let \(\mathrm{T}(u)\) be the trace of the endomorphism \(m_u: x \mapsto ux\) of the free A-module E (II, §4, no. 3). Then the mapping \(s\) defined by \(s(u) = \mathrm{T}(u)\) . \(e - u\) is an automorphism of the algebra E and \(s^2(u) = u\) for all \(u \in \mathbf{E}\) .

Let \((e, i)\) be a basis of E of type \((\alpha, \beta)\) ; then \(T(e) = 2\) , when \(s(e) = e\) , and \(T(i) = \beta\) , whence \(s(i) = \beta e - i\) . Hence \((e, s(i))\) is a basis of E, whose type is given by (3) with \(\gamma = \beta\) and \(\delta = -1\) , which again gives \((\alpha, \beta)\) ; it follows that \(s\) is an automorphism of the algebra E. As \(m_{s(u)} = sm_u s^{-1}\) , the endomorphisms \(m_u\) and \(m_{s(u)}\) of the A-module E have the same trace (II, §4, no. 3, Proposition 3), whence

\[
s ^ {2} (u) = \mathrm{T} (u). e - s (u) = \mathrm{T} (u). e - (\mathrm{T} (u). e - u) = u
\]

for all \(u \in E\) .

The automorphism \(s\) is called conjugation of the A-algebra \(\mathbf{E}\) and \(s(u)\) the conjugate of \(u\) .

If \(u = \xi e + \eta i\) , with \(\xi, \eta\) in \(A\) , then \(s(u) = (\xi + \beta \eta)e - \eta i\) , whence

\begin{enumerate}
\def\labelenumi{(\arabic{enumi})}
\setcounter{enumi}{3}
\tightlist
\item
\end{enumerate}

\[
\mathrm{T} (u) e = u + s (u) = (2 \xi + \beta \eta) e\tag{5}
\]

\[
u. s (u) = (\xi^ {2} + \beta \xi \eta - \alpha \eta^ {2}) e = N (u) e
\]

where we have written \(\mathrm{N}(u)=\xi^{2}+\beta\xi\eta-\alpha\eta^{2}\) . The elements \(\mathrm{T}(u)\) and \(\mathrm{N}(u)\) (or, when A and Ae are canonically identified, the elements \(\mathrm{T}(u)e\) and \(\mathrm{N}(u)e\) ) are called respectively the trace and norm of u.

When \(\beta = 0\) , the above formulae are simplified to

\[
s (\xi e + \eta i) = \xi e - \eta i, \quad \mathrm{T} (\xi e + \eta i) = 2 \xi , \quad \mathrm{N} (\xi e + \eta i) = \xi^ {2} - \alpha \eta^ {2}. \tag {6}
\]

Clearly T is a linear form on E and N is a quadratic form on E (IX, § 3, no. 4). As E is commutative and associative, it follows from (5) that

\[
\mathrm{N} (u v) = \mathrm{N} (u)) \mathrm{N} (v).\tag{7}
\]

For u to be invertible in E, it is necessary and sufficient that \(\mathrm{N}(u)\) be invertible in A. For, as \(\mathrm{N}(e)=1\) , the necessity of the condition follows from (7), writing \(v=u^{-1}\) . Conversely, if \(\mathrm{N}(u)\) is invertible in A, it follows from (5) that u is invertible and that

\[
u ^ {- 1} = (\mathrm{N} (u)) ^ {- 1} s (u).\tag{8}
\]

It can be proved that \(\mathbf{N}(u)\) is the determinant (§ 8, no. 1) of the endomorphism \(m_u\) (cf.~§ 9 no. 3, Example 1).

The following proposition gives the structure of quadratic algebras over a commutative field:

PROPOSITION 3. Let E be a quadratic A-algebra of type \((\alpha, \beta)\) .

\begin{enumerate}
\def\labelenumi{(\roman{enumi})}
\item
  If \(\mathbf{A}\) is a field and contains no element \(\zeta\) such that \(\zeta^2 = \alpha + \beta \zeta\) , \(\mathbf{E}\) is \(a\) (commutative) field (cf.~V, § 3).
\item
  If the ring A contains an element \(\zeta\) such that \(\zeta^2 = \alpha + \beta\zeta\) and \(\beta - 2\zeta\) is invertible (resp. zero), E is isomorphic to \(\mathbf{A} \times \mathbf{A}\) (resp. is of type \((0,0)\) ).
\end{enumerate}

We prove (i). Let \(\xi, \eta\) be two elements of A and \(u = \xi e + \eta i\) . If \(\eta \neq 0\) and we write \(\theta = -\xi \eta^{-1}\) , then \(\mathrm{N}(u) = \eta^2 (\theta^2 - \beta \theta - \alpha)\) by (5), whence \(\mathrm{N}(u) \neq 0\) by virtue of the hypothesis on A; if \(\eta = 0\) , then \(\mathrm{N}(u) = \xi^2\) . In any case, if \(u \neq 0\) , then \(\mathrm{N}(u) \neq 0\) , hence \(\mathrm{N}(u)\) is invertible in A and therefore \(u\) is invertible in E.

We now prove (ii). The canonical basis \((e_{1}, e_{2})\) of the algebra \(A \times A\) is of type \((0, 1)\) . We have seen (formula (3)) that E is of type

\[
(\alpha \delta^ {2} - \gamma^ {2} - \beta \gamma \delta , 2 \gamma + \beta \delta)
\]

for all \(\gamma \in \mathbf{A}\) and all \(\delta\) invertible in \(\mathbf{A}\) . If \(\beta - 2\zeta\) is invertible, take \(\delta = (\beta - 2\zeta)^{-1}\) and \(\gamma = -\zeta(\beta - 2\zeta)^{-1}\) ; then \(2\gamma + \beta\delta = 1\) and

\[
\alpha \delta^ {2} - \gamma^ {2} - \beta \gamma \delta = \delta^ {2} (\alpha - \zeta^ {2} + \beta \zeta) = 0;
\]

thus E is of type \((0,1)\) and hence isomorphic to A × A. If \(\beta - 2\zeta = 0\) , it has already been remarked that E is of type \((\alpha + \zeta^{2}, 0)\) and hence of type \((0,0)\) since \(\alpha + \zeta^{2} = 2\zeta^{2} - \beta\zeta = 0\) .

A quadratic A-algebra of type \((0,0)\) is also called an algebra of dual numbers over A.

